\documentclass[aspectratio=169,11pt,t]{beamer}

\usepackage{amsmath,amssymb,booktabs,array}
\usepackage[T1]{fontenc}
\usepackage{lmodern}

\definecolor{Ink}{HTML}{111111}
\definecolor{Muted}{HTML}{60646C}
\definecolor{Panel}{HTML}{F0F2F4}
\definecolor{Rule}{HTML}{B8BCC4}
\definecolor{Accent}{HTML}{3D8DFF}
\definecolor{AccentLight}{HTML}{DDEEFF}

\setbeamercolor{normal text}{fg=Ink,bg=white}
\setbeamercolor{frametitle}{fg=Ink,bg=white}
\setbeamercolor{structure}{fg=Accent}
\setbeamercolor{block title}{fg=Ink,bg=AccentLight}
\setbeamercolor{block body}{fg=Ink,bg=Panel}
\setbeamercolor{alerted text}{fg=Accent}

\setbeamerfont{title}{size=\fontsize{28}{32},series=\bfseries}
\setbeamerfont{subtitle}{size=\fontsize{15}{19}}
\setbeamerfont{frametitle}{size=\fontsize{16}{25},series=\bfseries}
\setbeamerfont{normal text}{size=\fontsize{13}{17}}
\setbeamerfont{block title}{size=\fontsize{13}{16},series=\bfseries}
\setbeamerfont{block body}{size=\fontsize{12}{16}}

\setbeamertemplate{navigation symbols}{}
\setbeamertemplate{itemize item}{\color{Accent}\raisebox{0.15ex}{\small$\bullet$}}
\setbeamertemplate{itemize subitem}{\color{Accent}\raisebox{0.15ex}{\scriptsize$\bullet$}}
\setbeamertemplate{enumerate item}{\color{Accent}\insertenumlabel.}
\setbeamertemplate{frametitle}{%
  \vspace{0.35cm}\insertframetitle\par
  \vspace{0.08cm}{\color{Rule}\rule{\textwidth}{0.5pt}}
}
\setbeamertemplate{footline}{%
  \leavevmode
  \hbox{%
    \begin{beamercolorbox}[wd=.86\paperwidth,ht=2.7ex,dp=1.2ex,leftskip=.65cm]{normal text}%
      \scriptsize\color{Muted}Energy-based basket intelligence
    \end{beamercolorbox}%
    \begin{beamercolorbox}[wd=.14\paperwidth,ht=2.7ex,dp=1.2ex,rightskip=.65cm plus1fil]{normal text}%
      \scriptsize\color{Muted}\hfill\insertframenumber
    \end{beamercolorbox}%
  }
}

\newcommand{\blue}[1]{\textcolor{Accent}{#1}}
\newcommand{\smallsource}[1]{\vfill{\tiny\color{Muted}#1}}

\title{Energy-based basket intelligence for retail}
\subtitle{One probability model for recommendation, generation and decision support}
\author{}
\date{}

\begin{document}

% 1
\begin{frame}[plain]
  \vspace{1.2cm}
  {\color{Accent}\large RETAIL BASKET MODEL}\par
  \vspace{0.8cm}
  {\usebeamerfont{title}\inserttitle\par}
  \vspace{0.45cm}
  {\usebeamerfont{subtitle}\color{Muted}\insertsubtitle\par}
  \vfill
  {\color{Rule}\rule{0.46\textwidth}{0.8pt}}\par
  \vspace{0.22cm}
  {\small\color{Muted}Version 4 model and current empirical interpretation}
\end{frame}

% 2
\begin{frame}{The model turns transactions into basket intelligence}
\begin{columns}[T,onlytextwidth]
  \column{0.31\textwidth}
  \begin{block}{Understand}
  Learn what different customers buy, which products appear together, and how basket
  composition changes across contexts.
  \end{block}

  \column{0.31\textwidth}
  \begin{block}{Predict}
  Complete a basket, rank candidate products, and generate plausible baskets for a known
  shopping context.
  \end{block}

  \column{0.31\textwidth}
  \begin{block}{Plan}
  Screen bundles, prices and promotions before spending money on controlled retailer
  experiments.
  \end{block}
\end{columns}

\vspace{0.22cm}
\begin{center}
\large\blue{Current position: predictive basket intelligence, not autonomous pricing.}
\end{center}
\end{frame}

% 3
\begin{frame}{The retailer receives five practical capabilities}
\begin{columns}[T,onlytextwidth]
  \column{0.48\textwidth}
  \begin{enumerate}
    \item \textbf{Basket recommendation}\par
          Rank products that best complete a customer's current basket.
    \item \textbf{Bundle discovery}\par
          Find supported product relationships beyond popularity alone.
    \item \textbf{Customer understanding}\par
          Summarize preferences, basket patterns and apparent price sensitivity.
  \end{enumerate}

  \column{0.48\textwidth}
  \begin{enumerate}
    \setcounter{enumi}{3}
    \item \textbf{Scenario screening}\par
          Compare predicted basket response under candidate prices or promotions.
    \item \textbf{Basket simulation}\par
          Generate baskets for testing, planning and model diagnostics.
  \end{enumerate}

  \vspace{0.12cm}
  \blue{Business use: reduce thousands of possible actions to a small,
  evidence-based shortlist.}
\end{columns}
\end{frame}

% 4
\begin{frame}{The problem is to model a complete basket, not isolated products}
\small
For a known shopping context $x$, we want one normalized probability for every feasible
non-empty basket $S$:

\[
\boxed{\quad p(S\mid x,\text{purchase trip})\quad}
\]

\begin{columns}[T,onlytextwidth]
  \column{0.48\textwidth}
  \begin{block}{Inputs}
  Customer, prices, promotions, calendar, store, available catalogue and declared history.
  \end{block}

  \column{0.48\textwidth}
  \begin{block}{Output}
  A probability distribution over sets of distinct products, with basket sizes from 1 to
  120.
  \end{block}
\end{columns}

\small Current data conditions on a purchase trip; predicting a visit requires a separate model.
\end{frame}

% 5
\begin{frame}{Every basket receives an energy, then energies become probabilities}
\[
p_{\Theta}(S\mid x)
=\frac{\exp\{E_{\Theta}(S;x)\}}{Z_{+}(x)},
\qquad
Z_{+}(x)=\sum_{T\in\Omega_x^{+}}\exp\{E_{\Theta}(T;x)\}.
\]

\begin{columns}[T,onlytextwidth]
  \column{0.48\textwidth}
  \begin{itemize}
    \item $S$: the basket being scored
    \item $x$: the observed shopping context
    \item $E_{\Theta}(S;x)$: basket energy or compatibility score
    \item $\Theta$: all learned model parameters
  \end{itemize}

  \column{0.48\textwidth}
  \begin{itemize}
    \item $\mathcal A_x$: products offered in context $x$
    \item $\Omega_x^{+}$: all non-empty subsets of $\mathcal A_x$ up to size 120
    \item $Z_{+}(x)$: the normalizer over every feasible basket
    \item Higher energy means higher probability after normalization
  \end{itemize}
\end{columns}
\end{frame}

% 6
\begin{frame}{The energy combines product value, interaction and basket shape}
\[
\boxed{
E_{\Theta}(S;x)=
\sum_{j\in S}b_j(x)
+\sum_{j<k;\,j,k\in S}\phi_j^{\top}\phi_k
-\sum_c\rho_c\binom{n_c(S)}{2}
-\rho_0(|S|)
}
\]

\begin{columns}[T,onlytextwidth]
  \column{0.24\textwidth}
  \textbf{$b_j(x)$}\par
  Contextual value of product $j$.

  \column{0.24\textwidth}
  \textbf{$\phi_j^{\top}\phi_k$}\par
  Low-rank relationship between products $j$ and $k$.

  \column{0.24\textwidth}
  \textbf{$\rho_c,n_c(S)$}\par
  Category effect and number of basket products in category $c$.

  \column{0.24\textwidth}
  \textbf{$\rho_0(|S|)$}\par
  Flexible preference over total basket size.
\end{columns}

\vspace{0.35cm}
$j,k$ index products, $c$ indexes affinity groups, and $|S|$ is the number of distinct
products in the basket.
\end{frame}

% 7
\begin{frame}{Contextual utility explains why the same product differs by trip}
\small
\[
\begin{aligned}
b_{jht}={}&\lambda_j+\theta_h^{\top}\alpha_j
-g_{hj}\,\Delta\log p_{jt}
+w_j^{\mathrm{dsp}}D_{jt}+w_j^{\mathrm{mlr}}M_{jt}\\
&+\mu_j^{\top}\delta_{w(t)}
+\zeta_j^{\top}\xi_{s(t)}
+\psi_j^{\top}r_{jht}.
\end{aligned}
\]

\begin{columns}[T,onlytextwidth]
  \column{0.48\textwidth}
  \begin{itemize}
    \item $j,h,t$: product, household and shopping time
    \item $\lambda_j$: product's baseline appeal
    \item $\theta_h^{\top}\alpha_j$: household-product preference
    \item $g_{hj}\ge0$: household-product price sensitivity
    \item $\Delta\log p_{jt}$: change in log price
  \end{itemize}

  \column{0.48\textwidth}
  \begin{itemize}
    \item $D_{jt},M_{jt}$: display and mailer promotion indicators
    \item $w_j^{\mathrm{dsp}},w_j^{\mathrm{mlr}}$: promotion responses
    \item $\mu_j^{\top}\delta_{w(t)}$: calendar effect
    \item $\zeta_j^{\top}\xi_{s(t)}$: store effect
    \item $\psi_j^{\top}r_{jht}$: declared history features
  \end{itemize}
\end{columns}
\end{frame}

% 8
\begin{frame}{The only fundamental computational obstacle is the normalizer}
\small
The observed basket is easy to score. The difficulty is the denominator:

\[
Z_{+}(x)=\sum_{T\in\Omega_x^{+}}\exp\{E_{\Theta}(T;x)\}.
\]

\begin{center}
\begin{minipage}{0.82\textwidth}
\begin{block}{Why direct enumeration fails}
With 5,455 products, the number of possible subsets is on the scale of $2^{5455}$.
Training, incidence probabilities and calibrated counterfactuals all require $Z_{+}$ or
its derivatives.
\end{block}
\end{minipage}
\end{center}

\end{frame}

% 9
\begin{frame}{The tractability theorem removes exponential catalogue growth}
\small
Let $z\sim\mathcal N(0,I_r)$, where $r$ is the interaction rank. Define

\[
w_j(z;x)=\exp\!\left\{b_j(x)-\tfrac12\|\phi_j\|^2+z^{\top}\phi_j\right\}.
\]

The Hubbard-Stratonovich and elementary-polynomial result is

\[
\boxed{
Z_{+}(x)=
\mathbb E_z\!\left[
\sum_{n=1}^{n_{\max}}e^{-\rho_0(n)}A_n(z,x)
\right].
}
\]

$A_n(z,x)$ is the exact total conditional weight of all size-$n$ baskets, computed by
elementary symmetric polynomials and category convolution.

\begin{center}
\blue{Result: an exact dynamic program plus an $r$-dimensional Gaussian integral.}
\end{center}
\end{frame}

% 10
\begin{frame}{Three simple identities make the theorem work}
\begin{enumerate}
  \item \textbf{Rewrite all pair interactions as one squared low-rank sum}
  \[
  \sum_{j<k\in S}\phi_j^{\top}\phi_k
  =\tfrac12\left\|\sum_{j\in S}\phi_j\right\|^2
  -\tfrac12\sum_{j\in S}\|\phi_j\|^2.
  \]

  \item \textbf{Use a Gaussian expectation for the squared term}
  \[
  \mathbb E_z[e^{z^{\top}v}]=e^{\|v\|^2/2}.
  \]

  \item \textbf{Use elementary symmetric polynomials to sum distinct products}

  The coefficient of degree $n$ collects every size-$n$ basket exactly once, including
  its category and total-size terms.
\end{enumerate}
\end{frame}

% 11
\begin{frame}{Training and every application use the same joint law}
For an observed basket $S_x$, training maximizes

\[
\ell_x(\Theta)=E_{\Theta}(S_x;x)-\log Z_{+}(x).
\]

\begin{columns}[T,onlytextwidth]
  \column{0.31\textwidth}
  \begin{block}{Observed signal}
  Raise the energy of baskets that actually occurred.
  \end{block}

  \column{0.31\textwidth}
  \begin{block}{Model correction}
  Lower probability assigned to alternatives through $\log Z_{+}$.
  \end{block}

  \column{0.31\textwidth}
  \begin{block}{One objective}
  Recommendation is not trained separately; it is a conditional query of the fitted law.
  \end{block}
\end{columns}

\end{frame}

% 12
\begin{frame}{Recommendation is exact because the common normalizer cancels}
\small
Let $R$ be the revealed basket and let product $j\notin R$ be a candidate completion.
Define its add-one score

\[
s_j(R,x)=E(R\cup\{j\};x)-E(R;x).
\]

Then

\[
\boxed{
P(j\text{ completes }R\mid x,\text{one item missing})
=\frac{e^{s_j(R,x)}}{\sum_{k\notin R}e^{s_k(R,x)}}.
}
\]

The score combines standalone utility, compatibility with products already in $R$,
category crowding and the basket-size increment. Ranking needs no partition-function
approximation.
\end{frame}

% 13
\begin{frame}{Basket generation follows the probability law in four stages}
\small
\begin{columns}[T,onlytextwidth]
  \column{0.48\textwidth}
  \begin{enumerate}
    \item \textbf{Draw the latent interaction state $z$}\par
          This summarizes the low-rank basket interaction direction.
    \item \textbf{Draw total size $n$}\par
          Use weights $e^{-\rho_0(n)}A_n(z,x)$ for $n=1,\ldots,120$.
  \end{enumerate}

  \column{0.48\textwidth}
  \begin{enumerate}
    \setcounter{enumi}{2}
    \item \textbf{Draw category counts}\par
          Allocate the selected total across affinity groups.
    \item \textbf{Draw distinct products}\par
          Sample the required number within each group using the conditional weights
          $w_j(z;x)$.
  \end{enumerate}
\end{columns}

\vspace{0.10cm}
\blue{Current output: distinct-product baskets. Units, margin and inventory need a
separate quantity model.}
\end{frame}

% 14
\begin{frame}{Incidence and price scenarios are derivatives or reweighted expectations}
\small
\begin{columns}[T,onlytextwidth]
  \column{0.48\textwidth}
  \begin{block}{Product incidence}
  \[
  \pi_j(x)=P(j\in S\mid x)
  =\frac{\partial\log Z_{+}(x)}{\partial b_j}.
  \]
  This includes utility, interactions, basket size and competition for basket positions.
  \end{block}

  \column{0.48\textwidth}
  \begin{block}{Price or promotion scenario}
  If action $a$ changes basket energy by
  $\Delta_a(S)=E_a(S)-E_{\mathrm f}(S)$, then
  \[
  \mathbb E_a[f(S)]
  =\frac{\mathbb E_{\mathrm f}[f(S)e^{\Delta_a(S)}]}
  {\mathbb E_{\mathrm f}[e^{\Delta_a(S)}]}.
  \]
  \end{block}
\end{columns}

\vspace{0.05cm}
\blue{Predictive response; causal only with credible identified price variation.}
\end{frame}

% 15
\begin{frame}{Interactions and segments turn probabilities into retailer insight}
\small
The identifiable interaction object is the Gram matrix
$K=\Phi\Phi^{\top}$, where $K_{jk}=\phi_j^{\top}\phi_k$.

\begin{columns}[T,onlytextwidth]
  \column{0.48\textwidth}
  \begin{block}{Product relationships}
  $K_{jk}$ measures how product $j$ changes product $k$'s add-one energy, holding the rest
  of the basket fixed. Category and size effects must also be included.
  \end{block}

  \column{0.48\textwidth}
  \begin{block}{Customer segments}
  Group household preference surfaces and summarize their products, basket patterns and
  apparent price response. Segments describe behaviour; they are not permanent identities.
  \end{block}
\end{columns}

\vspace{0.05cm}
\blue{Predictive bundle evidence, not proof of causal complementarity.}
\end{frame}

% 16
\begin{frame}{Current evidence supports basket intelligence, with a generation warning}
\small
\begin{columns}[T,onlytextwidth]
  \column{0.48\textwidth}
  \textbf{What is established}
  \begin{itemize}
    \item Gain over exact additive parent:
          $0.03275\pm0.00239$ nats per basket.
    \item Gains over Bernoulli, DPP and NDPP:
          $2.252$, $2.262$ and $1.812$ nats.
    \item Locked test MRR: $0.09525\pm0.00607$.
    \item Strong interaction-pair lift: $1.216$;
          matched controls: $0.998$.
  \end{itemize}

  \column{0.48\textwidth}
  \textbf{What is not yet calibrated}
  \begin{itemize}
    \item Generated mean size: $7.34$ versus observed $10.03$.
    \item Generated size variance: $63.83$ versus observed $136.28$.
    \item Interaction-only recommendation gain is not yet conclusive.
    \item Price responses are predictive, not causal revenue estimates.
  \end{itemize}
\end{columns}

\smallsource{Source: locked Version 4 household-size rank-one evaluation.}
\end{frame}

% 17
\begin{frame}{Retail deployment should move from prediction to controlled learning}
\begin{columns}[T,onlytextwidth]
  \column{0.31\textwidth}
  \begin{block}{1. Understand}
  Train on historical baskets. Deliver segments, recommendations, supported product pairs
  and scenario rankings.
  \end{block}

  \column{0.31\textwidth}
  \begin{block}{2. Test}
  Randomize a small set of model-supported bundles or promotions. Record exposure,
  non-purchase, quantity and margin.
  \end{block}

  \column{0.31\textwidth}
  \begin{block}{3. Optimize}
  Calibrate incremental profit, add inventory and sequential state, then optimize a
  budget-constrained policy.
  \end{block}
\end{columns}

\vspace{0.08cm}
\begin{center}
\small\blue{RL can optimize a calibrated sequential environment; it cannot repair a biased one.}
\end{center}
\end{frame}

% 18
\begin{frame}{The immediate value is a smaller, better decision space}
\vspace{0.35cm}
\begin{center}
\begin{minipage}{0.86\textwidth}
{\Large\bfseries The model gives the retailer one coherent view of basket behaviour.}\par
\vspace{0.5cm}
{\large It connects customer context, product value, interactions, basket size and price
response in a normalized probability law.}\par
\vspace{0.6cm}
{\large\color{Accent}Use it now to recommend, discover and shortlist.\par}
\vspace{0.22cm}
{\large\color{Accent}Use controlled experiments to turn those predictions into decisions.\par}
\end{minipage}
\end{center}
\vfill
{\color{Rule}\rule{\textwidth}{0.8pt}}
\end{frame}

\end{document}
